\documentclass[12pt, a4paper]{article}

% ─── Paquetes ────────────────────────────────────────────────────────────────
\usepackage[utf8]{inputenc}
\usepackage[T1]{fontenc}
\usepackage[spanish]{babel}
\usepackage{geometry}
\geometry{margin=2.5cm}
\usepackage{graphicx}
\usepackage{hyperref}
\usepackage{enumitem}
\usepackage{titlesec}
\usepackage{parskip}
\usepackage{booktabs}
\usepackage{array}
\usepackage{xcolor}
\usepackage{fancyhdr}
\usepackage{caption}

% ─── Cabecera / Pie ──────────────────────────────────────────────────────────
\pagestyle{fancy}
\fancyhf{}
\rhead{Realidad Virtual y Aumentada}
\lhead{EHU/UPV}
\cfoot{\thepage}

% ─── Metadatos del documento ─────────────────────────────────────────────────
\hypersetup{
    colorlinks=true,
    linkcolor=blue!60!black,
    urlcolor=blue!60!black,
    pdftitle={Memoria Proyecto RVA - DeliveryVR},
    pdfauthor={Iker Hidalgo, Adriana Bernal}
}

% ─────────────────────────────────────────────────────────────────────────────
\begin{document}

% ─── Portada ─────────────────────────────────────────────────────────────────
\begin{titlepage}
    \centering
    \vspace*{2cm}
    \makebox[\textwidth]{
        \includegraphics[height=1.75cm]{EHU_Logo.pdf}
        \hfill
        \includegraphics[height=1.75cm]{KISa_Logo.jpg}
    }
    \\[3cm]

    {\Huge \textbf{DeliveryVR}}\\[0.8em]
    {\LARGE Realidad Virtual y Aumentada}\\[3cm]

    \begin{tabular}{rl}
        \textbf{Autores:} & Iker Hidalgo \\
                          & Adriana Bernal \\[1em]
        \textbf{Curso:}   & 2025--2026 \\
    \end{tabular}

    \vfill
    {\large \today}
\end{titlepage}

% ─── Índice ──────────────────────────────────────────────────────────────────
\tableofcontents
\newpage

% ─────────────────────────────────────────────────────────────────────────────
\section{Introducción}

Este documento describe el desarrollo de \textbf{DeliveryVR}, un videojuego de realidad virtual creado como proyecto para la asignatura de Realidad Virtual y Aumentada.

El objetivo principal del proyecto ha sido diseñar una experiencia inmersiva orientada a dispositivos de realidad virtual, explorando mecánicas basadas en la interacción mediante controladores y movimientos físicos del usuario. Para ello, se ha implementado un prototipo arcade en el que el jugador debe entregar paquetes depositándolos en diferentes zonas de entrega señalizadas, montado en un camión en movimiento. 

Este     proyecto ha sido desarrollado íntegramente desde cero utilizando Godot y OpenXR, sin partir de plantillas jugables ni proyectos base preconfigurados. Tanto la lógica de interacción en realidad virtual como los sistemas de puntuación, movimiento del vehículo, manipulación de objetos y construcción del escenario han sido implementados específicamente para este proyecto.

Con el objetivo de facilitar la consulta y reproducibilidad del proyecto, el código fuente completo y los recursos utilizados se encuentran disponibles públicamente en el siguiente repositorio de GitHub:

\url{https://github.com/Hidalgo7/Realidad_Virtual_XR}

% ─────────────────────────────────────────────────────────────────────────────
\section{Descripción del proyecto}

En DeliveryVR, el jugador asume el rol de un repartidor que viaja sobre la plataforma trasera de un camión mientras este circula de forma autónoma por un entorno urbano. Durante el trayecto existen varias zonas de entrega señalizadas en el suelo mediante indicadores visuales iluminados en color verde. El objetivo consiste en lanzar y encestar las cajas en dichas zonas antes de que el vehículo continúe avanzando y las deje atrás. 

\subsection{Jugabilidad}

En la siguiente lista se describen los elementos principales de jugabilidad implementados.

\begin{itemize}[leftmargin=*]
    \item \textbf{Generación de cajas}: Al pulsar el \textit{trigger} del mando derecho se generan cajas que el jugador deberá utilizar para completar los repartos en las diferentes zonas señalizadas.
    \item \textbf{Agarre}: Tras soltar una caja generada (soltar el \textit{trigger}), el jugador puede volver a recogerla si se encuentra a una distancia adecuada mediante los botones \textit{grip} de ambos controladores.
    \item \textbf{Lanzamiento}: Al lanzar una caja esta se libera con una velocidad y dirección acorde a la velocidad del control que la estaba sujetando.
    \item \textbf{Puntuación}: Cada caja que entra en una zona de entrega suma 100 puntos al marcador global, y todas las zonas requieren completar un número aleatorio de entregas entre 1 y 5.
    \item \textbf{Movimiento del vehículo}: El camión se desplaza por la ciudad siguiendo una trayectoria predefinida.
\end{itemize}

\subsection{Entorno}

El escenario consiste en una ciudad de estilo \textit{low-poly}, compuesta por distintos edificios, carreteras con curvas y tramos rectos, y diversos elementos urbanos distribuidos a lo largo del recorrido. Gran parte de los recursos visuales empleados para construir este entorno proceden del paquete de assets \textit{City Builder Bits} \cite{citybuilderbits}, los cuales han sido adaptados e integrados específicamente para las necesidades del proyecto.


La simplicidad visual del entorno favorece un rendimiento estable en dispositivos de realidad virtual y evita sobrecargar al usuario con un exceso de estímulos visuales, reduciendo así el riesgo de mareo por movimiento. A pesar de mantener una estética sencilla, el recorrido incorpora múltiples escenas y elementos urbanos distribuidos dinámicamente, como controles de policía, paradas de taxi, fuentes, señales de tráfico y mobiliario urbano. Estos elementos no cumplen únicamente una función decorativa, sino que también contribuyen a reforzar la percepción de movimiento, profundidad y presencia espacial del usuario, aumentando la sensación de inmersión dentro del entorno virtual sin comprometer el rendimiento general del sistema.


Durante el desarrollo se priorizó la comodidad del usuario y la estabilidad visual del entorno. Para ello se tomaron varias decisiones específicas de diseño:

\begin{itemize}[leftmargin=*]
    \item El jugador tiene capacidad de movimiento y orientación libres, aunque el espacio de acción está limitado únicamente al maletero del camión.
    \item El recorrido del vehículo mantiene velocidades suaves y giros amplios para minimizar la desorientación.
    \item El estilo \textit{low-poly} reduce la carga visual y favorece un rendimiento estable en el visor.
    \item Las zonas objetivo utilizan colores brillantes y animaciones pulsantes para mejorar la legibilidad espacial.
    \item La interacción se basa en movimientos naturales de las manos mediante seguimiento XR.
\end{itemize}

Estas decisiones siguen principios habituales de diseño de experiencias inmersivas y buscan reducir la fatiga visual mientras se mantiene una sensación de presencia convincente.

% ─────────────────────────────────────────────────────────────────────────────
\section{Arquitectura y estructura del proyecto}

\subsection{Organización de ficheros}

El proyecto se organiza en tres directorios principales:

\begin{itemize}
    \item \texttt{assets}. Recursos 3D para construir el entorno: edificios, coches, elementos de mobiliario urbano y el vehículo del jugador.
    \item \texttt{scripts}. Aloja los scripts de lógica: gestor de partida, controladores de ambas
    manos, ruta del camión, detección de cajas, etc.
    \item \texttt{scenes}. Contiene las escenas del juego, que combinan los recursos visuales con la lógica de jugabilidad. Entre ellas están la escena principal, el camión, las zonas de puntuación, las cajas y los controladores de la manos virtuales.
\end{itemize}

\subsection{Lógica principal}

La funcionalidad del prototipo se concentra en varios scripts. Las
escenas definen la estructura de nodos, modelos y colisiones, mientras que los scripts coordinan
la interacción XR, el movimiento del camión, la puntuación y la retroalimentación visual.

\begin{itemize}[leftmargin=*]

    \item \textbf{Control de manos}: En los scripts \texttt{right\_hand\_controller.gd} y \texttt{left\_hand\_controller.gd} se concentra la interacción principal con los controladores. Extienden de la clase \texttt{XRController3D} y gestionan las señales de pulsación y liberación de los botones trigger y grip. 
    
    Para el caso concreto de la mano derecha al recibir la señal \texttt{trigger\_click} se instancia la escena de la caja añadiendola a la escena actual, busca dentro de ella el elemento \texttt{RigidBody3D} y llama a la función \texttt{\_pick\_up\_object} del nodo \texttt{FunctionPickup} de Godot XR Tools, lo que hace que la caja quede sujeta a la mano.

    Adicionalmente, los scripts de ambas manos implementan una búsqueda de objetos cercanos con la función \texttt{\_find\_pickable()} que se activa mediante la señal \texttt{grip\_click}. Dicha función recorre los hijos de la escena actual filtrando los nodos de tipo \texttt{XRToolsPickable} y selecciona el más cercano al marcador \texttt{GrabDetectPoint} dentro de un radio razonable para agarrar la caja. Al soltar el botón, se llama a la utilidad \texttt{drop\_object()} de \texttt{FunctionPickup}, soltando la caja de forma natural.

     \item \textbf{Movimiento del camión}: Está gestionado por el script \texttt{path\_follow\_3d.gd}. El script extiende \texttt{PathFollow3D} y, en cada ciclo del cálculo de físicas controlado por la función \texttt{\_physics\_process}, incrementa \texttt{progress} usando la velocidad exportada. La escena principal usa este nodo sobre un \texttt{Path3D} y sincroniza el cuerpo animable del camión mediante \texttt{RemoteTransform3D}.

    \item \textbf{Áreas de entrega}: La lógica interna de las zonas de entrega se gestiona mediante el fichero \texttt{score\_area.gd}.
    
    En la función inicial \texttt{\_ready()} se asigna aleatoriamente el número máximo de cajas en un rango del 1 al 5 necesarias mediante la función \texttt{randi\_range()}. Además se actualiza el texto que indica el progreso de la zona y conecta la señal \texttt{body\_entered} que detecta las cajas que entrarán en la zona. Cuando entra un cuerpo del grupo \texttt{cajas}, incrementa el contador, suma los puntos configurados, elimina la caja y, si se alcanza el objetivo, marca la zona como completada desactivando el flag \texttt{monitoring}.

    \item \textbf{Puntuación}: Se gestiona mediante \texttt{GameManager.gd} que funciona como \textit{autoload} global y centraliza la puntuación de la partida. Mantiene la variable \texttt{score}, busca dinámicamente el \texttt{Label3D} del visor y actualiza el texto del marcador cada vez que una zona invoca \texttt{add\_score(points)}. De este modo, la zona de entrega no necesita conocer la estructura completa de la interfaz XR.

\end{itemize}

% ─────────────────────────────────────────────────────────────────────────────
\section{Resultados y conclusiones}

El resultado final es una experiencia de realidad virtual funcional que cumple los objetivos
planteados: el jugador puede generar cajas y repartirlas lanzándolas a diferentes zonas mientras el camión circula por la ciudad.

Desde el punto de vista académico, el proyecto ha permitido aplicar de forma práctica
conceptos clave de la asignatura:

\begin{itemize}[leftmargin=*]
    \item \textbf{Presencia e inmersión}: La combinación de manos virtuales representadas
    fielmente y la física realista del lanzamiento contribuyen a la sensación de estar
    realmente sobre el camión.
    \item \textbf{Confort en VR}: El entorno limitado de acción y el movimiento suave del desplazamiento del camión busca minimizar la cinetosis, una de las principales barreras para la adopción de la realidad virtual.
    \item \textbf{Estándares de la industria}: La integración con OpenXR y las herramientas de
    Godot XR Tools reflejan el uso de estándares y bibliotecas de referencia actuales en el
    desarrollo de aplicaciones XR.
\end{itemize}

Como líneas de trabajo futuro se podrían contemplar: añadir un contador de tiempo visible
en el visor, introducir cajas de diferentes tamaños con distinto valor de puntuación, incorporar
efectos de sonido espacial al lanzamiento e impacto, o ampliar el recorrido con nuevas secciones
de ciudad y obstáculos dinámicos.

% ─────────────────────────────────────────────────────────────────────────────
\begin{thebibliography}{9}

\bibitem{godotxr}
Godot Engine Documentation.
\textit{XR in Godot}.\\
\url{https://docs.godotengine.org/en/stable/tutorials/xr/index.html}

\bibitem{godotxrtools}
Godot XR Tools.
\textit{GitHub Repository}.\\
\url{https://github.com/GodotVR/godot-xr-tools}

\bibitem{citybuilderbits}
Kay Lousberg.
\textit{City Builder Bits}.\\
\url{https://kaylousberg.itch.io/city-builder-bits}

\end{thebibliography}

\end{document}
